% ============================================================================
\clearpage
\section{Datos y reconstrucción: la primera capa de robustez}\label{sec:datos}
% ============================================================================

\ficha{Cotizaciones bid/ask de cadenas completas, subyacente sincronizado, tasas, dividendos y eventos
corporativos, con sellos de tiempo de mercado y de recepción.}
{Filtrar, estimar forward y descuento, convertir opciones americanas, ajustar una superficie sin
arbitraje y extraer la densidad.}
{Distribución $\Q_t$ a plazo constante, cuantiles, bandas de sensibilidad y soporte identificado.}
{Nada se sobrescribe; densidad no negativa, masa próxima a uno y media igual al forward; colas sin
soporte marcadas.}

Todo lo que viene después depende de esta etapa. Un error aquí no produce una señal débil: produce una
señal falsa con apariencia de precisión. Por eso la reconstrucción se diseña como una cadena de pasos
auditables (\figref{fig:flujo_q}), cada uno con un control explícito.

\begin{figure}[H]
\centering
\resizebox{\linewidth}{!}{% Reconstrucción de Q: de cotizaciones a cuantiles con bandas.
\begin{tikzpicture}[
  n/.style={caja=qazul,text width=3.15cm,minimum height=15mm},
  chk/.style={nota,text width=3.15cm,align=center,text=tinta2},
  node distance=6mm and 5.5mm]
\node[n] (s1) {\textbf{1. Snapshot inmutable}\\ bid, ask, tamaños, sellos de tiempo, subyacente};
\node[n,right=of s1] (s2) {\textbf{2. Filtros de calidad}\\ cruces, bid nulo, spreads, datos rancios};
\node[n,right=of s2] (s3) {\textbf{3. Forward y descuento}\\ paridad, tasas y dividendos};
\node[n,right=of s3] (s4) {\textbf{4. Desamericanización}\\ equivalente europeo o exclusión};
\node[n,below=12mm of s4] (s5) {\textbf{5. Ajuste en precios}\\ SSVI restringido, pérdida por banda};
\node[n,left=of s5] (s6) {\textbf{6. Contraste}\\ ajuste monótono y convexo};
\node[n,left=of s6] (s7) {\textbf{7. Densidad}\\ Breeden--Litzenberger y controles};
\node[n,left=of s7] (s8) {\textbf{8. Cuantiles y bandas}\\ plazo constante, colas marcadas};
\foreach \a/\b in {s1/s2,s2/s3,s3/s4,s5/s6,s6/s7,s7/s8} \draw[flecha] (\a) -- (\b);
\draw[flecha] (s4) -- (s5);
\node[chk,below=0.5mm of s1] {sin sobrescribir nada};
\node[chk,below=0.5mm of s2] {motivo de cada exclusión};
\node[chk,below=0.5mm of s3] {forward coherente con paridad};
\node[chk,above=0.5mm of s4.north] {error de conversión acotado};
\node[chk,below=0.5mm of s5] {restricciones por construcción};
\node[chk,below=0.5mm of s6] {discrepancia = riesgo de modelo};
\node[chk,below=0.5mm of s7] {$q\ge 0$, masa $\approx 1$, media $=F$};
\node[chk,below=0.5mm of s8] {sensibilidad, no confianza};
\end{tikzpicture}
}
\caption{Cadena de reconstrucción de $\Q$. Debajo de cada paso, el control que debe superar antes de
alimentar al siguiente.}
\label{fig:flujo_q}
\end{figure}

\subsection{Snapshots inmutables y datos puntuales en el tiempo}\label{sec:snapshots}

Cada captura se guarda como un \emph{snapshot inmutable}: nunca se corrige en sitio; una corrección es
un registro nuevo con su propia hora de publicación. Esto permite responder, para cualquier instante
pasado, la única pregunta válida en un backtest: \emph{¿qué se sabía en ese momento?}
La \tabref{tab:campos} resume los campos mínimos y por qué importan.

\begin{table}[H]
\centering\small
\caption{Campos mínimos de cada snapshot.}
\label{tab:campos}
\begin{tabularx}{\linewidth}{@{}l Y@{}}
\toprule
\encab{Campo} & \encab{Por qué importa}\\
\midrule
Bid, ask y tamaños & La banda $[b,a]$ es la observación; el mid es solo un punto interior.\\
Hora de mercado y de recepción & Separan cuándo ocurre un dato de cuándo se conoce (\figref{fig:puntual}).\\
Subyacente sincronizado & Un desfase de segundos contamina el forward y la volatilidad implícita.\\
Strike, vencimiento, estilo, multiplicador & Definen el contrato; el estilo americano exige conversión.\\
Dividendos, tasas y eventos corporativos & Determinan forward, descuento y ajustes de contrato.\\
Revisiones y open interest & Se guardan con su hora efectiva de disponibilidad, no con la del dato.\\
\bottomrule
\end{tabularx}
\end{table}

\begin{figure}[H]
\centering
\resizebox{0.96\linewidth}{!}{% Datos puntuales en el tiempo: cuándo ocurre un dato frente a cuándo se conoce.
\begin{tikzpicture}[x=2.05cm,y=1cm]
% Región posterior al corte
\fill[qazul!6] (3.75,2.75) rectangle (6.3,-0.25);
% Pistas: hora de mercado (arriba) y hora de recepción (abajo); 06:00 -> x = 0
\draw[guia,line width=0.7pt] (0,2.2) -- (6.3,2.2);
\draw[guia,line width=0.7pt] (0,0) -- (6.3,0);
\node[titulo,anchor=east,align=right] at (-0.08,2.2) {ocurre\\[-1pt] {\mdseries\tiny(hora de mercado)}};
\node[titulo,anchor=east,align=right] at (-0.08,0) {se conoce\\[-1pt] {\mdseries\tiny(hora de recepción)}};
\foreach \h/\e in {0/06:00,1/07:00,2/08:00,3/09:00,4/10:00,5/11:00,6/12:00}{
  \draw[guia] (\h,2.2) -- (\h,2.28); \draw[guia] (\h,0) -- (\h,-0.08);
  \node[nota,anchor=north] at (\h,-0.1) {\e};}
% Corte de decisión
\draw[qazul,line width=0.9pt] (3.75,2.95) -- (3.75,-0.45);
\node[nota,anchor=south,text=qazuloscuro] at (3.75,2.95) {\textbf{corte de decisión: 09:45 ET}};
\node[nota,anchor=north west,text=qazuloscuro] at (3.8,-0.45) {recibido después del corte: no existe para esa decisión};
% Open interest del cierre anterior, publicado antes de la apertura
\draw[flecha,draw=tenue] (0.08,2.2) -- (1.0,0);
\node[nota,anchor=south west] at (0.0,2.3) {open interest del cierre de $t-1$};
% Fuente con 15 minutos de retraso: dato de 09:30 recibido a las 09:45
\draw[flecha,draw=qnaranja,line width=0.8pt] (3.5,2.2) -- (3.74,0);
\node[nota,anchor=east,align=right] at (3.52,1.05) {fuente con 15 min de retraso:\\ dato de 09:30 recibido 09:45};
% Cotización en tiempo real
\draw[flecha,draw=qazul,line width=0.8pt] (3.71,2.2) -- (3.715,0);
\node[nota,anchor=south east] at (3.7,2.3) {cotización 09:44:58};
% Revisión posterior de un dato anterior
\draw[flecha,draw=critico,line width=0.8pt] (3.6,2.2) -- (5.08,0);
\node[nota,anchor=west,align=left] at (4.55,1.35) {revisión publicada a las 11:05\\ de un dato de las 09:36};
\end{tikzpicture}
}
\caption{Hora en que ocurre un dato frente a hora en que se conoce. En el corte de las 09:45 solo son
utilizables las flechas que llegan a la línea inferior antes del corte: la cotización en tiempo real,
el open interest del cierre anterior y, con una fuente retrasada, lo ocurrido a las 09:30. La revisión
publicada a las 11:05 no existe para esa decisión aunque se refiera a un dato anterior.}
\label{fig:puntual}
\end{figure}

\subsection{Contrato de datos y fuentes candidatas}\label{sec:contrato}

Antes de adquirir historia extensa conviene validar una muestra contra un contrato explícito. El
primer contrato debe demostrar, como mínimo:

\begin{itemize}
\item cadenas \emph{completas}, incluidos contratos ya vencidos, y paginación sin huecos;
\item cobertura histórica \emph{a la hora de decisión}, no solo al cierre;
\item acceso a cotizaciones de compra y venta. Las barras OHLC de operaciones no sustituyen una cadena
contemporánea de cotizaciones: una operación puede ser antigua, aislada o fuera del spread;
\item sellos de tiempo suficientes para reconstruir la sincronía con el subyacente.
\end{itemize}

\begin{table}[H]
\centering\small
\caption{Fuentes candidatas según la documentación citada en la propuesta. Cobertura, licencia y precio
deben cotizarse; no se asumen.}
\label{tab:fuentes}
\begin{tabularx}{\linewidth}{@{}l Y Y@{}}
\toprule
\encab{Fuente} & \encab{Lo que documenta} & \encab{Uso previsto}\\
\midrule
Alpaca & Históricos de opciones desde febrero de 2024; el feed indicativo modifica cotizaciones
\citep{alpaca}. & Explorar la integración. No se asume suficiente para estudiar cambios diarios pequeños
de distribución.\\
Cboe DataShop & Snapshots NBBO por intervalos con historia desde enero de 2012; entrega intradía estándar
con 15 minutos de retraso \citep{cboe}. & Investigación histórica. Una decisión en apertura exige una
fuente con latencia adecuada o retrasar explícitamente la decisión.\\
\bottomrule
\end{tabularx}
\end{table}

\begin{criterio}
Cadenas reconstruibles en el instante de decisión, con cobertura y costo conocidos, validadas sobre una
muestra antes de comprar historia extensa.
\end{criterio}

\subsection{Cadena cruda: controles, paridad y asimetría observada}\label{sec:cadena_cruda}

Las primeras preguntas del plan de trabajo (\texttt{docs/plan\_de\_trabajo.md}) se responden con la cadena
\emph{cruda}: si el encarecimiento relativo de calls frente a puts, o el residuo de paridad del mismo
strike, anticipa algo. Estas medidas se calculan antes de cualquier ajuste y se guardan como una serie
propia: una superficie sin arbitraje impone la paridad y borraría justamente el residuo que se quiere
estudiar. El módulo \texttt{quantileflow/cadenas.py} implementa esta sección. Su clase \texttt{Captura} es
la versión mínima del contrato de datos: una fila por actualización de cada contrato, con bid, ask,
tamaños y sello de tiempo, más el subyacente sincronizado, la tasa, los dividendos, el estilo de
ejercicio, la hora de corte y la fecha. Sus arreglos son de solo lectura.

\paragraph{Controles por fila} De cada contrato cuenta solo la última actualización anterior o igual al
corte; si esa actualización falla un control, el contrato queda fuera y no se recurre a una cotización más
vieja. Cada fila excluida conserva \emph{todos} sus motivos (\tabref{tab:motivos}), de modo que el informe
diario dice qué se excluyó y por qué. Los umbrales son provisionales hasta auditar una muestra real en la
fase~0. La captura recibe además alertas si el sello del subyacente se aleja del corte o cae en la
ventana de apertura, cuando el valor de un índice puede incluir componentes que aún no han abierto.

\begin{table}[H]
\centering\small
\caption{Motivos de exclusión de una fila de la cadena cruda.}
\label{tab:motivos}
\begin{tabularx}{\linewidth}{@{}l Y@{}}
\toprule
\encab{Motivo} & \encab{Condición y razón}\\
\midrule
\texttt{posterior\_al\_corte} & Sello posterior al corte: no se conocía al decidir.\\
\texttt{reemplazada} & Hay una actualización posterior del mismo contrato antes del corte.\\
\texttt{antes\_de\_apertura} & Cotización de preapertura o de la sesión anterior.\\
\texttt{ventana\_de\_apertura} & Primeros minutos de la sesión (cinco por omisión): cotizaciones inestables.\\
\texttt{desfasada} & Más antigua que la edad máxima al corte (60~s): no refleja el subyacente actual.\\
\texttt{sin\_bid} & Bid nulo: el precio solo queda acotado por arriba.\\
\texttt{cruzada}, \texttt{bloqueada} & Bid mayor que ask, o igual: dato inconsistente o de fuentes
desincronizadas.\\
\texttt{sin\_tamano} & Tamaño menor que el mínimo en bid o en ask.\\
\texttt{spread\_ancho} & Spread mayor que el 50\% del mid.\\
\texttt{fuera\_de\_cotas} & Viola cotas sin modelo: call por encima del subyacente, put por encima del
strike o, si es americana, ask menor que el ejercicio inmediato.\\
\bottomrule
\end{tabularx}
\end{table}

\paragraph{Residuo de paridad fuera de muestra} Para cada par $j$ de call y put válidos del mismo strike,
\begin{equation}
r_j=(C_j-P_j)-D_{-j}\,(F_{-j}-K_j),\qquad
\bigl[\underline r_j,\ \overline r_j\bigr]=\bigl[C^{b}_j-P^{a}_j,\ C^{a}_j-P^{b}_j\bigr]-D_{-j}\,(F_{-j}-K_j),
\label{eq:residuo_paridad}
\end{equation}
donde $C_j$ y $P_j$ son mids, los superíndices $b$ y $a$ indican bid y ask, y $(F_{-j},D_{-j})$ se estiman
por mínimos cuadrados ponderados de $C-P$ sobre $K$ con los \emph{demás} pares, con pesos $1/\ell^2$ y
$\ell_j=(C^a_j-C^b_j)+(P^a_j-P^b_j)$ el ancho de la banda. Ninguna cotización evalúa su propia paridad, de
modo que la medida no es circular. Un par sale de la estimación del forward de los demás si su residuo
supera tres medios anchos y cinco desviaciones robustas del resto, pero conserva su residuo. El residuo se
reporta en unidades de precio y en múltiplos del ancho: si $|r_j|\le\ell_j/2$, la banda contiene el cero y
la diferencia no supera el costo de negociar, así que no se etiqueta como señal. La
\figref{fig:paridad_desfasadas} muestra por qué el control de edad es imprescindible: tres puts de hace
cinco minutos, con el subyacente un 1\% más bajo, dan residuos de $-0.6$, $-1.9$ y $-3.2$ anchos que,
sin ese control, pasarían por señales; los demás pares no se contaminan.


\paragraph{Dividendos y ejercicio anticipado} El forward implícito incorpora los dividendos sin necesidad
de declararlos. Un forward contractual que omite un dividendo en efectivo fabrica, en todos los strikes, un
residuo igual a su valor presente ($-1.2$ en las pruebas, con un dividendo de 1.2). Con opciones americanas
la paridad deja de ser una igualdad: cada cotización se desamericaniza con el árbol binomial y el modelo de
dividendo \emph{escrowed} (\secref{sec:americanas}), valorando la americana y la europea en el mismo
árbol, y el residuo solo se interpreta si supera la banda \emph{más} las primas de ejercicio estimadas y
el subyacente está sincronizado. Tratar como europeas unas americanas con dividendo sesga el forward
implícito (0.2 sobre 98.8 en las pruebas) aunque ningún residuo salga de su banda.

\grafica[0.62]{s2_paridad_desfasadas}{Residuo de paridad del mismo strike, en múltiplos del ancho bid--ask,
con el forward estimado sin el propio par. Los círculos son los pares que superan los controles; los
triángulos, tres puts desfasados que solo aparecerían si no se comprobara la edad de la cotización. Entre
las líneas discontinuas, la diferencia cabe en el spread.}{fig:paridad_desfasadas}

\paragraph{Asimetría call/put} Un call a $+3$\% y un put a $-3$\% no forman un par de paridad: su
comparación mide asimetría de precios de cola. Se usan strikes simétricos en logaritmo, $Fe^{\pm d}$ con
$d=\ln 1.03$, y se compara primero en volatilidad implícita, $\sigma_P(-d)-\sigma_C(d)$, con una banda que
sale de las volatilidades de bid y ask. Las primas brutas no deben compararse con la unidad: con una
sonrisa simétrica en $k$,
\begin{equation}
P\bigl(Fe^{-d}\bigr)=e^{-d}\,C\bigl(Fe^{d}\bigr),\qquad
\varrho_d=\frac{P(Fe^{-d})}{e^{-d}\,C(Fe^{d})}-1,
\label{eq:simetria_pc}
\end{equation}
de modo que $\varrho_d=0$ indica ausencia de asimetría. Como una distancia fija cambia de significado con
el nivel de volatilidad y el plazo, la comparación se repite con deltas comparables,
$\sigma_P(\Delta=-0.25)-\sigma_C(\Delta=0.25)$, usando la delta forward de cada opción. La pendiente local
$\partial\sigma/\partial k$ en $k=0$ se estima por mínimos cuadrados ponderados con $|k|\le5$\%
(\figref{fig:asimetria_sonrisa}). Solo se interpola entre strikes válidos separados a lo sumo 2\% en
$k$; si el tramo necesario tiene un hueco mayor, la medida queda «no identificada».

\grafica[0.62]{s2_asimetria_sonrisa}{Volatilidades implícitas de bid y ask de las cotizaciones OTM válidas y
las tres medidas de asimetría: strikes log-simétricos en $\pm\ln 1.03$, deltas de 25 y pendiente local. En
este mundo sintético, la diferencia en $\pm3$\% es de 5.8 puntos, frente a 5.76 en la sonrisa
verdadera.}{fig:asimetria_sonrisa}

\paragraph{Series separadas y cambios diarios} La fila diaria del informe lleva las medidas observadas
(\texttt{obs\_*}) y, si se ajusta una superficie, las que se leen de ella (\texttt{aj\_*}), nunca
mezcladas: en el ajuste, el residuo de paridad es cero por construcción. Los cambios respecto de la
captura anterior registran días naturales, sesiones transcurridas y tipo de comparación
(apertura--apertura o cierre--apertura); un fin de semana cuenta como tres días y una sesión, no como una
sesión ordinaria.

\begin{noconcluir}
Las pruebas de esta sección usan cadenas sintéticas con cotizaciones desfasadas, dividendos, strikes
ausentes y ventanas de apertura. Verifican que el procesamiento separa efectos de los datos de una posible
señal; no demuestran que el residuo o la asimetría anticipen rendimientos, ni la calidad de datos reales.
\end{noconcluir}

\subsection{De cotizaciones a precios ajustados}\label{sec:ajuste_precios}

\paragraph{Forward y descuento} Para opciones europeas, la paridad put--call
$C(K)-P(K)=D\,(F-K)$ permite estimar $D$ y $F$ por regresión ponderada de $C-P$ sobre $K$
(\secref{sec:cadena_cruda}). Con opciones americanas con dividendos la paridad es solo una desigualdad: se
aplica a los precios desamericanizados, con la prima de ejercicio estimada como parte de la
incertidumbre, y se contrasta con el forward que dan la curva de tasas y los dividendos esperados.

\paragraph{Selección y filtros} Se trabaja con opciones fuera del dinero (puts bajo el forward, calls por
encima), más líquidas y con menor valor intrínseco. Se excluyen, registrando el motivo, las cotizaciones
cruzadas, rancias o con spread anómalo (\tabref{tab:motivos}). Una cotización con bid nulo no se descarta sin más: su ask sigue
siendo una cota superior útil, aunque no identifica la cola (\figref{fig:cotizaciones}).

\grafica{s2_cotizaciones_precio}{Cotizaciones OTM a 30 días en el mundo sintético. Cada segmento vertical
es un intervalo bid--ask; los triángulos son cotizaciones con bid nulo, que solo acotan el precio por
arriba. La curva es el ajuste SSVI en espacio de precios.}{fig:cotizaciones}

\paragraph{Ajuste en espacio de precios} El midpoint no es una verdad exacta: cualquier precio dentro de
$[b_i,a_i]$ es compatible con la cotización. Por eso el residuo de cada contrato es su \emph{distancia a la
banda}, medida en medios spreads $s_i=\max\{(a_i-b_i)/2,\ \text{tick}\}$:
\begin{equation}
\min_{\vartheta\in\Theta}\ \sum_i \omega_i\left[\frac{(b_i-C_i(\vartheta))^+ + (C_i(\vartheta)-a_i)^+}{s_i}\right]^2
+\eta\sum_i \omega_i\left[\frac{C_i(\vartheta)-m_i}{s_i}\right]^2+\mathcal{R}(\vartheta),
\label{eq:perdida_banda}
\end{equation}
donde $C_i(\vartheta)$ es el precio del modelo, $m_i$ el mid, $\eta$ un peso pequeño que solo desempata
entre precios igualmente compatibles y $\mathcal R$ una regularización, por ejemplo hacia el ajuste de la
sesión anterior. Ajustar en precios y no en volatilidades evita dar peso excesivo a las alas, donde un
spread de un tick equivale a varios puntos de volatilidad (\figref{fig:cotizaciones_vol}).

\grafica{s2_cotizaciones_vol}{La misma cadena en volatilidad implícita. Los intervalos se ensanchan en las
alas y las cotizaciones con bid nulo solo dan cotas superiores cada vez más altas: en volatilidad, esas
cotas parecen ``información'' sobre la cola derecha que en realidad no existe.}{fig:cotizaciones_vol}

\grafica[0.6]{s2_perdida_banda}{Pérdida de una cotización con bid 1.00 y ask 1.20. La pérdida cuadrática
al mid penaliza precios perfectamente compatibles; la distancia a la banda es casi nula dentro del
intervalo y crece fuera de él.}{fig:perdida_banda}

\subsection{SSVI y ausencia de arbitraje estático}\label{sec:ssvi}

El modelo base es SSVI \citep{gatheral2014}, que describe la varianza implícita total $w=\sigma^2T$ como
función del log-moneyness $k=\ln(K/F)$ y de la varianza total en el dinero $\theta_T$ de cada vencimiento:
\begin{equation}
w(k,\theta_T)=\frac{\theta_T}{2}\Bigl\{1+\rho\,\varphi(\theta_T)\,k+\sqrt{\bigl(\varphi(\theta_T)k+\rho\bigr)^2+1-\rho^2}\Bigr\},
\qquad \varphi(\theta)=\frac{\eta}{\theta^{\gamma}(1+\theta)^{1-\gamma}}.
\label{eq:ssvi}
\end{equation}
El parámetro $\rho$ controla la asimetría, $\varphi$ la curvatura y el peso de las alas, y $\theta_T$ el
nivel por vencimiento. Una SVI sin restricciones no garantiza ausencia de arbitraje; SSVI admite
condiciones suficientes sencillas:

\begin{itemize}
\item \textbf{Calendario:} $\theta_T$ no decreciente en $T$ y
$0\le\partial_\theta\bigl(\theta\varphi(\theta)\bigr)\le\rho^{-2}\bigl(1+\sqrt{1-\rho^2}\bigr)\varphi(\theta)$.
\item \textbf{Mariposa:} $\theta\varphi(\theta)(1+|\rho|)<4$ y $\theta\varphi(\theta)^2(1+|\rho|)\le 4$.
\item \textbf{Función de potencia:} con $0<\gamma\le 1/2$ y $\eta(1+|\rho|)\le 2$ se cumplen ambas familias
de condiciones, siempre que $\theta_T$ sea no decreciente.
\end{itemize}

La primera condición de mariposa es coherente con la fórmula de momentos de \citet{lee2004}: la pendiente
asintótica de la varianza total no puede superar 2.

\begin{notatecnica}[Restricciones por construcción]
En la implementación de referencia las restricciones no se verifican a posteriori: se imponen con una
reparametrización sin restricciones, $\rho=0.999\tanh(a)$, $\gamma=\tfrac12\,\sigma(b)$,
$\eta=2\sigma(c)/(1+|\rho|)$ y $\theta_i=\sum_{j\le i}e^{d_j}$, con $\sigma$ la función logística. Cuando
el ajuste queda en la frontera de la región suficiente (como en la \figref{fig:ssvi_region}), conviene
registrarlo: la región es suficiente, no necesaria, y una restricción activa puede estar sesgando la
forma.
\end{notatecnica}

\grafica[0.6]{s2_ssvi_rebanadas}{Rebanadas de varianza total ajustadas para 7, 30, 60 y 91 días en el mundo
sintético. No se cruzan: esa es la forma gráfica de la ausencia de arbitraje de calendario.}{fig:ssvi_rebanadas}

\grafica[0.6]{s2_ssvi_region}{Región suficiente de la función de potencia en el plano $(|\rho|,\eta)$ con
$\gamma\le 1/2$. El cuadrado es el valor verdadero del mundo sintético; el círculo, el ajuste, que queda
en la frontera.}{fig:ssvi_region}

La prueba directa de mariposa usa la función $g$ de Durrleman, en la forma de \citet{gatheral2014}:
con $w'$ y $w''$ derivadas respecto de $k$,
\begin{equation}
g(k)=\Bigl(1-\frac{k\,w'}{2w}\Bigr)^2-\frac{(w')^2}{4}\Bigl(\frac1w+\frac14\Bigr)+\frac{w''}{2},
\label{eq:g}
\end{equation}
y la densidad de $x=\ln(S_T/F)$ es $p(k)=g(k)\,(2\pi w)^{-1/2}\exp\{-d_-(k)^2/2\}$ con
$d_-(k)=-k/\sqrt{w}-\sqrt{w}/2$. Una rebanada tiene densidad no negativa si y solo si $g\ge 0$
(\figref{fig:funcion_g}).

\grafica[0.6]{s2_funcion_g}{Función $g(k)$ para una rebanada SSVI restringida y para una SVI sin
restricciones con parámetros publicados como contraejemplo en \citet{gatheral2014}. Donde $g<0$ la
densidad implícita sería negativa: un arbitraje de mariposa.}{fig:funcion_g}

\paragraph{Comparación independiente} Como contraste se ajusta, sin forma paramétrica de volatilidad, una
densidad discreta $q_j\ge 0$ sobre una malla $x_j$:
\begin{equation}
C(K)=D\sum_j q_j\,(x_j-K)^+,\qquad \sum_j q_j=1,\qquad \sum_j q_j x_j=F.
\label{eq:convexo}
\end{equation}
Cualquier $q\ge0$ produce precios decrecientes y convexos en $K$, y las dos igualdades fijan la masa y la
media del forward. Con un precio auxiliar $p_i\in[b_i,a_i]$ por cotización, la pérdida por banda se vuelve
un problema de mínimos cuadrados lineales con cotas, convexo y rápido de resolver, al que se añade una
penalización de suavidad sobre $q$. Donde SSVI y este ajuste discrepan hay \emph{riesgo de modelo}, no un
dato \citep{aitsahalia1998}.

\subsection{Opciones americanas y equivalentes europeos}\label{sec:americanas}

SPY y las acciones individuales tienen opciones americanas; la identidad de extracción de la
\secref{sec:bl} es europea. Hay dos caminos: convertir cada cotización en un equivalente europeo con un
modelo que incluya ejercicio anticipado y dividendos, o excluir las observaciones donde la conversión sea
demasiado incierta. La conversión estándar (\emph{desamericanización}) tiene dos pasos: invertir el precio
americano con un árbol binomial \citep{cox1979} o diferencias finitas con dividendos discretos para obtener
una volatilidad, y valorar con ella la opción europea equivalente \citep{burkovska2018}.

Seleccionar opciones fuera del dinero reduce el problema pero no lo elimina
(figuras~\ref{fig:americanas_puts} y~\ref{fig:americanas_calls}). El error de conversión depende del
modelo y de los dividendos supuestos; por eso se contrasta dentro de los escenarios de incertidumbre
(\secref{sec:bandas}) en lugar de tratarlo como exacto.

\grafica[0.6]{s2_americanas_puts}{Sesgo, en puntos de volatilidad, de tratar puts americanas como europeas
(líneas continuas), comparado con medio spread expresado en volatilidad (discontinuas). A 91 días el sesgo
de puts cerca del dinero ya es comparable con medio spread, aunque estén fuera del dinero.}{fig:americanas_puts}

\grafica[0.6]{s2_americanas_calls}{Lo mismo para calls con un dividendo en efectivo dentro del plazo. Las
calls dentro del dinero antes de la fecha ex-dividendo tienen una prima de ejercicio anticipado grande; las
calls fuera del dinero apenas se ven afectadas.}{fig:americanas_calls}

\subsection{Extracción de la densidad}\label{sec:bl}

Para calls europeas y descuento determinista $D(t,T)$ se cumple \citep{breeden1978}
\begin{equation}
q^{\Q}_t(K,T)=\frac{1}{D(t,T)}\,\frac{\partial^2 C_t(K,T)}{\partial K^2},\qquad
\Q_t(S_T\le K)=1+\frac{1}{D(t,T)}\,\frac{\partial C_t(K,T)}{\partial K}.
\label{eq:bl}
\end{equation}
La identidad se aplica al \emph{precio ajustado}, nunca a diferencias de cotizaciones crudas: una segunda
derivada amplifica el ruido de redondeo y de spread hasta producir densidades negativas
(figuras~\ref{fig:densidad_cruda} y~\ref{fig:densidad_ajustada}).

\grafica[0.62]{s2_densidad_cruda}{Segundas diferencias de los mids crudos (mundo sintético con strikes cada
punto). La curva discontinua es la densidad verdadera, desconocida en la práctica; los puntos rojos son
valores negativos, imposibles para una densidad.}{fig:densidad_cruda}

\grafica[0.62]{s2_densidad_ajustada}{La misma identidad sobre precios ajustados. La densidad SSVI integra
0.9998 en la malla y su media es $1.0001F$; el ajuste monótono--convexo, independiente, coincide en el
cuerpo.}{fig:densidad_ajustada}

Los controles mínimos de cada densidad son tres, y se registran con su valor numérico:

\begin{enumerate}
\item \textbf{No negatividad} en toda la malla evaluada.
\item \textbf{Masa próxima a uno}. La masa que falta está en colas sin strikes y se reporta; no se
redistribuye en silencio.
\item \textbf{Primer momento compatible con el forward}: $\E^{\Q}[S_T]=F$.
\end{enumerate}

También se documentan el error numérico (paso de la malla, diferencias finitas) y la sensibilidad a la
extrapolación de strikes extremos.

\paragraph{Cuantiles sin masa inventada} Los cuantiles se obtienen invirtiendo una CDF \emph{anclada}: la
masa a la izquierda de la malla sale de la pendiente del precio (segunda identidad de~\eqref{eq:bl}), no
de integrar la densidad desde un extremo desconocido. No se normaliza la masa ni se recortan valores
negativos. Cada nivel $u$ se devuelve con su estado: «identificada» si cae entre $\Q(X\le x_{\min})$ y
$\Q(X\le x_{\max})$, con $[x_{\min},x_{\max}]$ el tramo respaldado por strikes fiables, y «no
identificada», sin valor, en otro caso. Una densidad negativa invalida todos los niveles. Normalizar
habría sido engañoso: a una normal a la que le falta el 2.3\% de la masa izquierda, la normalización le
desplaza la mediana 0.029 desviaciones típicas y le atribuye una cola que nadie observó.

\begin{noconcluir}
La media de $\Q$ está fijada por el forward. No es una expectativa de rentabilidad real y no contiene, por
construcción, ninguna opinión del mercado sobre el signo del rendimiento esperado.
\end{noconcluir}

\subsection{Conjunto de superficies plausibles y colas no identificadas}\label{sec:bandas}

Un solo ajuste oculta cuánto podría variar la distribución sin dejar de ser compatible con las
cotizaciones. Por eso se mantiene un \emph{conjunto} de superficies plausibles, coherentes con los spreads
y las restricciones: se sortean precios objetivo dentro de cada $[b_i,a_i]$, se reajustan ambos métodos y
se repite con supuestos alternativos de dividendos y desamericanización. De ese conjunto salen bandas para
cuantiles y probabilidades de cola (figuras~\ref{fig:bandas_densidad} y~\ref{fig:bandas_cuantiles}).

\grafica[0.62]{s2_bandas_densidad}{Bandas de densidades compatibles con los spreads, en escala logarítmica.
La banda SSVI es estrecha incluso donde no hay strikes fiables (más allá de los cuadrados): esa precisión
la pone la forma paramétrica, no los datos. El ajuste monótono--convexo muestra cuánto se ignora
realmente en las colas.}{fig:bandas_densidad}

\grafica[0.62]{s2_bandas_cuantiles}{Ancho de la banda de sensibilidad de cada cuantil. Los cuadrados marcan
el rango identificado $[u_{\min},u_{\max}]$, definido por los strikes fiables más extremos; fuera de él
los cuantiles se reportan como no identificados.}{fig:bandas_cuantiles}

\begin{noconcluir}
Estas bandas son de \emph{sensibilidad}: describen qué distribuciones son compatibles con los datos y los
supuestos. No son intervalos de confianza estadísticos, porque no hay un modelo de observación que lo
justifique. Si falta soporte en una cola, se marca como no identificada; no se rellena con precisión
artificial.
\end{noconcluir}

\begin{criterio}
Coherencia matemática (controles de la \secref{sec:bl} superados en cada sesión) y estabilidad de
cuantiles y colas frente a perturbaciones plausibles de spreads y supuestos.
\end{criterio}
